\section{Geometric positive transfers and special complete intersections}

The abstract abelian-type transfer remains distinct from a transfer that
preserves the actual incidence form. The following statements give
such geometric transfers and identify the extra special-fibre directions
that they do not reach. Their positivity conclusions are conditional
on the unrefereed 032 input whenever that qualification is displayed.

\subsection*{C44-W06. A finite abelian cover gives an explicit inherited polarization}

\begin{theorem}[C44-W06; conditional on unrefereed 032, Corollary 8.4]
Let $X$ be smooth projective of odd dimension $d=2m-1$ over an
algebraically closed field. Suppose that there is a finite surjective
morphism $f:A\to X$ from an abelian variety. For every rational ample
class $h$ on $X$, put $h_A=f^*h$ and $e=\deg f$. The primitive
incidence quotient $B_X=P^m(X)_{\rm prim}$ has the signed rational
polarization of cote 44, Th.~4. More precisely the induced map
\[
 i_f:B_X\longrightarrow B_A=P^m(A)_{\rm prim}
\]
has finite kernel, and its actual incidence classes satisfy
\begin{equation}\label{c44:finite-cover-polarization}
 i_f^*\ell_{B_A}=e\ell_{B_X},\qquad
 \ell_{B_X}=\frac{(-1)^m}{2}\Delta^*d_{B_X}.
\end{equation}
In particular this is a positive transfer on a concrete class of
nonabelian varieties of geometric abelian type. It does not use
Milne's additional transfer assertion from JAC-P20.
\end{theorem}

\begin{proof}
A finite morphism between smooth varieties of the same dimension is
flat; its degree includes inseparable multiplicity. Push--pull gives
$f_*f^*=e$, and $h_A$ is ample. Lefschetz inverses for $X$ are
algebraic: in degree $i\leq d$, an inverse to $L_h^{d-i}$ is
\[
 \frac1e f_* (L_{h_A}^{d-i})^{-1}f^*.
\]
Indeed $f^*$ commutes with the indicated Lefschetz operators and
the formula composes with $L_h^{d-i}$ to the identity. The inverse
on the abelian variety is algebraic. The same pull--push construction
with its Kunneth projectors supplies those of $X$. Thus the middle
primitive projector is algebraic and self-transpose. Saito's
incidence construction and transpose compatibility, used in
JAC-T4.1, give its nondegenerate factor $B_X$.

Pullback takes primitive cycle families on $X$ to primitive cycle
families on $A$, since $h_Af^*=f^*h$. It therefore induces $i_f$.
Pushforward induces a rational left inverse up to $e$, so $i_f$ has
finite kernel. For two normalized families $z,z'$ on parameter
varieties, projection formula gives the equality of divisorial
incidence correspondences
\[
 (p_{12})_*\bigl((1\times f)^*z\cdot(1\times f)^*z'\bigr)
   =e\,(p_{12})_*(z\cdot z').
\]
By the incidence functoriality of JAC-T4.1, this is the equality
$(i_f\times i_f)^*d_{B_A}=e d_{B_X}$. Diagonal pullback gives
\eqref{c44:finite-cover-polarization}, with the same sign and
normalization as JAC-T4.3. That theorem makes $\ell_{B_A}$ ample
conditional on 032. Pullback by a homomorphism with finite kernel
is ample on its source, and division by $e>0$ preserves rational
ampleness. This proves the conclusion. At the motive level
$e^{-1}f_*f^*=1$ exhibits $h(X)$ as an actual Chow-motive summand
of $h(A)$, so the stated abelian-type scope is also explicit.
\end{proof}

\paragraph{Concrete nonabelian examples.}
In characteristic different from two, take
$A=E_0\times E_1\times\cdots\times E_{d-1}$ with $d\geq3$ odd,
choose a nonzero two-torsion point $a\in E_0$, and quotient by
\[
 (x_0,x_1,x_2,\ldots)\longmapsto(x_0+a,-x_1,x_2,\ldots).
\]
The action is free because of its first coordinate. Its quotient
is smooth projective, the cover is finite etale, and the invariant
$H^1(\mathcal O)$ has dimension $d-1$, so the quotient is not an
abelian variety. In characteristic two use a nonzero three-torsion
translation on $E_0$ and an order-three automorphism of
$E_1:y^2+y=x^3$, acting by $x\mapsto\zeta_3x$. The same argument
works. These examples prove that C44-W06 has nonabelian content in
every characteristic.

\subsection*{C44-W07. The fixed part for every smooth complete intersection}

\begin{proposition}[C44-W07; fixed-part positivity conditional on unrefereed 032]
Let $A$ have dimension $N=2m+c-1$, let $h$ be ample, and let
$k:Y\hookrightarrow A$ be a smooth successive intersection of
$c$ divisors of class $h$, with $\dim Y=n=2m-1$. There is no
generic-last-section hypothesis in this statement. Define
\[
 r=(L_h^c)^{-1}k_*,\qquad
 F=k^*\ker L_h^{c+1},\qquad E=\ker k_*\subset H^n(Y),
\]
and let $e_A$ be the primitive projector on $H^n(A)$. Then
\[
 \Pi_F=k^*e_A r,\quad \Pi_E=1-k^*r,\qquad
 H^n(Y)_{\rm prim}=F\mathbin\perp E
\]
are algebraic self-adjoint projectors and an orthogonal decomposition.
They induce a rational isogeny decomposition of primitive incidence
quotients
\[
 P^m(Y)_{\rm prim}\sim B_F\times B_E.
\]
Its signed incidence class is nondegenerate on both factors and is
a rational polarization on $B_F$, conditional on 032. Consequently
the full Th.~4 assertion for this particular $Y$ is equivalent to
ampleness of its signed incidence class on $B_E$. In particular it
holds whenever $B_E=0$.
\end{proposition}

\begin{proof}
The fixed-and-variable projector proof of H9-T13 uses smooth weak
Lefschetz and the identity $[Y]=h^c$, before any pencil monodromy is
introduced. These hypotheses hold here. Explicitly,
$rk^*=1$, and $k^*r$ is the orthogonal projector onto the ambient
middle cohomology. The primitive condition for an ambient class
$k^*u$ is $L_h^{c+1}u=0$: the Gysin map on $H^{n+2}(Y)$ is
injective by dual weak Lefschetz. The same injectivity makes every
element of $\ker k_*$ primitive. The primitive projector of $A$
is algebraic; the projectors of $Y$ outside its middle degree come
from the weak-Lefschetz identifications with $A$ and their duals.
Thus the displayed middle projectors are actual algebraic
correspondences. Self-adjointness and nondegeneracy give the
incidence decomposition exactly as in JAC-T4.1.

For clarity, the positive scalar in the fixed-part test also remains
unchanged. If a normalized fixed family pulled to a pointed curve
$C$ of genus $b\geq1$ has class $\xi$, put $\eta=(1\otimes r)\xi$
and lift its parameter degree-one component to an algebraic class
$w$ on $J(C)\times A$. Then $h^{c+1}w=0$ and, for
$H=\theta+h$,
\[
 H^{b+c}w=0,\qquad
 \int_{J(C)\times A}w^2H^{b+c-1}
   =\frac{(b+c-1)!}{c!}\int_{C\times Y}\xi^2.
\]
This is H9-T13, equation (H9.47); its proof needs no genericity.
The coefficient is positive. Conditional on 032, the signed curve
degree is nonnegative, with equality only for a zero mixed class.
The nondegenerate signed divisor is therefore ample by the
nef-plus-nondegenerate criterion of JAC-T4.3. Orthogonality makes
ampleness of the full signed divisor equivalent to ampleness of
both factors, proving the final assertions.
\end{proof}

\subsection*{C44-W08. Specialization of the inherited incidence polarization}

\begin{theorem}[C44-W08; conditional on unrefereed 032]
Let $\mathscr Y\to U$ be the smooth locus of a genuine Lefschetz
pencil in the setting of H9-T13, and let $s\in U$ be any geometric
point. Write $B_{\bar\eta}=P^m(Y_{\bar\eta})_{\rm prim}$ for the
geometric generic primitive incidence quotient. After a finite
extension of the local generic field, $B_{\bar\eta}$ has good
reduction at $s$, and its special fibre maps with finite kernel to
$P^m(Y_s)_{\rm prim}$. The actual signed incidence form is positive
on this inherited image. Hence Th.~4 holds for $Y_s$ whenever
\begin{equation}\label{c44:no-incidence-jump}
 \dim P^m(Y_s)_{\rm prim}=\dim P^m(Y_{\bar\eta})_{\rm prim}.
\end{equation}
More particularly, if the generic variable incidence factor of the
pencil is nonzero, then Th.~4 holds for every smooth geometric fibre
of the pencil. If that generic factor is zero, only the new
variable incidence directions that occur at a special fibre remain
outside this theorem.
\end{theorem}

\begin{proof}
Work over the strict henselian local trait at $s$, or its completion,
and pass to a finite extension of its fraction field over which an
abelian representative $B$ and a normalized rational Poincare cycle
$z$ are defined. These are finite data; rational descent is the one
fixed in JAC-P00. As in H9-T5, equations (H9.10)--(H9.11), the
mixed class of $z$ injects the Tate module of $B$ into the appropriate
Tate twist of $H^n(Y_{\bar\eta})$. Indeed its self-incidence is
nondegenerate, so a vector killed by the family map would be in
the radical of that incidence pairing. Smooth proper base change
makes the target unramified at this smooth fibre. Thus $B$ has
good reduction by the Neron--Ogg--Shafarevich criterion, just as in
H9-T5. One exact primary author source for the criterion is Milne,
\emph{Complex Multiplication}, Theorem 6.12, p.~51, which cites
Serre--Tate (1968), Theorem 1.

Let $\mathscr B$ denote the resulting abelian scheme over the trait.
Take the closure of a rational-cycle representative of $z$ in
$\mathscr B\times\mathscr Y$. These schemes are smooth over a
regular trait, so the refined fibre restrictions and intersection
products are defined. Its special-fibre mixed class is the
specialization of the generic mixed class. In particular it is
still in parameter degree one and is primitive on $Y_s$; these
conditions are kernels of maps of the smooth-proper cohomology
systems. Rigidification or degree-one projection, if desired,
alters no incidence class.

Form the actual relative self-incidence divisor by the same proper
pushforward of the product of the two cycle families. Clear its
rational denominator. Its signed half-diagonal pullback is a line
bundle $\mathscr L$ on $\mathscr B$, and refined base change says
that its two fibre classes are exactly the corresponding signed
incidence classes, up to the one positive clearing factor. The
generic class is ample by H9-T13.

Here is the specialization step without an assertion that arbitrary
special cycles lift. Choose a relatively ample line bundle
$\mathscr H$ on $\mathscr B$. Flatness and the Hilbert polynomial
make
\[
 \chi(L_t+uH_t)=\frac{(c_1(L_t)+u c_1(H_t))^{\dim B}}{(\dim B)!}
\]
the same real polynomial in $u$ for the generic and special fibres
$t$. On the generic fibre this polynomial is positive for every
$u\geq0$, since $L_{\eta}+uH_{\eta}$ is ample. Thus on the
special fibre $\chi(L_s)\ne0$ and the polynomial has no positive
real root. Mumford's index theorem gives index zero for $L_s$;
the nondegenerate effective criterion makes it ample. These are
exactly the two abelian-variety criteria used in the checked proof
of JAC-T4.3 (Mumford, \emph{Abelian Varieties}, Chapter III,
Section 16, and Chapter II, Section 6, Application 1).

The specialized family induces a homomorphism
$u:B_s\to P^m(Y_s)_{\rm prim}$. Saito's transpose and incidence
compatibility identifies $u^*\ell_s$ with the ample class just
constructed. Therefore $u$ has finite kernel: restriction to any
positive-dimensional connected kernel would make an ample class
zero. This proves positivity on the inherited image. Under
\eqref{c44:no-incidence-jump} the map is an isogeny, and ampleness
descends through that isogeny.

For the final assertion apply the same construction to the generic
variable factor alone. If that factor is nonzero, the irreducibility
after finite-index restriction in H9-T13, or H9-T5 (H9.11), forces
its family map to fill the entire variable local system $E$.
Specialization preserves this rank. The special variable incidence
factor injects into $E_s$ by its own nondegenerate self-incidence,
so its dimension is at most $\operatorname{rank}(E)/2$. The
inherited image already has that dimension. It therefore maps
isogenously onto the special variable incidence factor and gives
its positive signed form. C44-W07 supplies the fixed factor for
every smooth fibre. This proves the stated all-smooth-fibre scope.
\end{proof}

\paragraph{The remaining special-fibre gate.}
When the generic variable incidence quotient vanishes, the
construction specializes only the fixed incidence quotient. A
special fibre may have a new variable incidence factor. Neither
smooth proper base change nor constancy of the rank of $E$ makes
that factor extend as algebraic cycle families through the pencil.
Its signed positivity is the smaller unresolved statement. Thus
C44-W08 does not turn H9-T13 into an unconditional assertion about
every smooth complete intersection; it proves precisely the
rank-preserving specializations and the nonzero-generic-variable
case, with 032 retained as the positivity input.

\subsection*{C44-W09. The full geometric-abelian-type extension with its extra input exposed}

\begin{proposition}[C44-W09; conditional on unrefereed 032 and the separate JAC-P20 transfer]
Assume, in addition to 032, Milne's Proposition 4.2 in
\emph{Grothendieck's standard conjecture of Lefschetz type over
finite fields}, version 14 July 2022, in the total-cohomology form
specified and audited in JAC-P20. Then JAC-T4.3 holds for every
smooth projective odd-dimensional variety of geometric abelian
type, with its actual signed primitive incidence divisor.
\end{proposition}

\begin{proof}
JAC-P16 gives algebraic Lefschetz and primitive projectors. For
every pointed test curve $C$, the product $J(C)\times X$ is still
of geometric abelian type. JAC-P20's actual abelian cohomology
injection and the additional transfer assertion give the Hodge
standard conjecture for this product. The construction of
JAC-T4.2, using the algebraic Abel inverse on the first factor,
now gives an algebraic primitive class $w$ on $J(C)\times X$ and
the exact identity
\[
 \int_{J(C)\times X}w^2(\theta+h)^{b-1}
       =(b-1)!\int_{C\times X}\xi^2.
\]
The Hodge-standard input on this particular product makes its
signed value nonnegative. Saito's nondegenerate incidence factor
and JAC-T4.3's curve criterion then make the signed divisor ample.
This argument proves the stated conditional corollary; it does
not establish the additional transfer assertion itself.
\end{proof}

\paragraph{Framework and precise dependency.}
C44-W06 proves a positive scalar comparison of the actual forms,
not merely an embedding of motives. C44-W07--W08 use smooth
proper base change, SGA 7's fixed and variable projectors, and
the extension of actual cycle correspondences across a regular
trait. Their target is cote 44, scans 008--009, Th.~4 and its
primitive incidence higher Picard variety. The new residual gate
is the positive sign on a special new variable incidence factor,
not the already checked generic pencil sign. For all varieties
of geometric abelian type, JAC-P20 still leaves the comparison
between an abstract transported positive form and the particular
geometric form as an additional dependency. C44-W09 records
the exact implication if that dependency is supplied.

